
This paper directly measured the statistical independence of SE\_N5 and min\_logprob as uncertainty signals for hallucination detection in open-weight LLMs on short-answer factual QA. On TriviaQA dev with Llama-3.1-8B-Instruct across N=2500 prompts, the Pearson correlation between these signals is |r| = 0.026 --- near-orthogonal, far below both the pre-registered independence threshold (0.70) and the reparameterization boundary (0.85). This result is not a marginal pass; it indicates that SE, which aggregates semantic entropy over N=5 stochastic samples, and min\_logprob, which extracts the minimum token confidence from a single greedy decode, operate as genuinely distinct uncertainty mechanisms.

The main contributions are: (1) Pearson |r|=0.026 confirming near-orthogonality at N=2500; (2) partial $R^{2}=0.0005$ (LRT p=0.442, LRT $chi^{2}=1.632$, df=2) showing that independence does not yield linear conditional predictive utility on this benchmark --- a powered null result for the pre-registered effect size, with the gate outcome (EXPLORE\_N10) disclosed transparently; (3) Spearman $\rho$(SE, judge)=$-0.026$ validating the cross-model LM-judge design as a low-circularity evaluation protocol; and (4) a fully reproducible, checkpoint-aware pipeline that completed the N=2500 evaluation.

Future work includes: nonlinear ensemble evaluation (h-e1-v2); extending to N=10 stochastic samples per prompt per the gate recommendation; cross-model transfer characterization (Llama-calibrated ensemble on Qwen-2.5-7B TruthfulQA); orthogonality zone analysis; bootstrapped CIs on primary metrics; and seed-stability verification. The findings motivate combining these signals through nonlinear methods rather than linear logistic regression. Careful statistical characterization of signal independence --- not just standalone AUROC comparison --- is the appropriate foundation for principled ensemble design in LLM hallucination detection.

